\hypertarget{appendix-6af5921bc853}{}\label{referencestart:appendix-6af5921bc853}
\section{Pointer arithmetic and provenance}\label{appendix-6af5921bc853-pointer-arithmetic-and-provenance}

\textbf{Curriculum depth:} Developed · \textbf{First seen:} C003 · \textbf{Latest development:} C008

\subsection{Definition and mechanism}\label{appendix-6af5921bc853-definition-and-mechanism}

M008 adds null/object/one-past distinctions, same-array representable subtraction and round-trip limits. Provenance proposals are research, not adopted C++23 rules; byte inspection differs from arbitrary typed access.

A scalar member behaves as a one-element array for pointer arithmetic, not as an array containing adjacent sibling members. \&x+1 can be a one-past pointer but cannot be dereferenced to access y. Actual array storage supports traversal; a span only views an already-valid range. Layout/size assertions do not authorize sibling traversal.

A pointer is a typed access path constrained by array bounds, alignment, and object lifetime, not merely an integer. Within an array, p+i advances i elements; one-past is permitted for traversal but not dereference. Byte views inspect representations; uintptr\_t, when provided, is an integer conversion facility, not permission to access arbitrary storage.

\subsection{Prerequisites}\label{appendix-6af5921bc853-prerequisites}

\begin{itemize}
\tightlist
\item
  \hyperref[appendix-18ee4f913c3c]{Addresses and virtual memory} --- Mappings give numeric locations their process-specific storage identity.
\item
  \hyperref[appendix-0b104bcba06b]{Storage duration and object lifetime} --- Access requires an object and backing storage that are still valid.
\end{itemize}

\subsection{Used by}\label{appendix-6af5921bc853-used-by}

\begin{itemize}
\item
  \href{../02-Mathematics-for-Visual-Computing/main.pdf\#nameddest=appendix-ffdc1d4e5350}{Matrix storage indexing and strides (Mathematics for Visual Computing, Reference notes: Matrix storage indexing and strides)}
\item
  \hyperref[appendix-3f7e15effba6]{Named members versus array elements} --- uses this prerequisite to establish its entity, lifetime, representation or lookup contract.
\item
  \hyperref[appendix-afae48c31800]{Null pointers and optional access}
\item
  \hyperref[appendix-f9b8b9188c7f]{Borrow invalidation and logical identity}
\item
  \href{../04-Light-Transport-and-Image-Synthesis/main.pdf\#nameddest=appendix-749ba51a6c25}{GPU vertex specification (Light Transport and Image Synthesis, Reference notes: GPU vertex specification)} \#\# Recall and next depth
\end{itemize}

Explain p+1 for int versus byte and why a numeric address is insufficient.

\textbf{Pending:} Full lifetime replacement, ownership/concurrency policies and cross-target ABI experiments; the opening examples do not establish those later guarantees.

\subsection{Source and navigation}\label{appendix-6af5921bc853-source-and-navigation}

\hyperref[appendix-fd090d982a15-m008---parts-113-126]{C008 development and corrections}

\hyperref[01-foundations-i-ch002-where-the-bytes-live]{C003 teaching chapter} · \hyperref[appendix-fd090d982a15-m003---parts-41-60]{Part definitions} · \hyperref[appendix-4ad47bdd269a]{Dependency Map} · \hyperref[01-foundations-i-ch003-bringing-the-object-model-back-to-bytes]{C007 development}
